\section{The ordinary line at a split irregular CM point}
\label{sec:cubic-irregular-saturation}

This section concerns the actual cohomological ordinary line, rather than
an abstract induced representation with the same traces.  The local
coefficient ring in the first two statements is
\(A=E[[X]]\), where \(E\) is a characteristic-zero field; its residue
field is denoted by \(k=E\).  All assertions about a ray tower below
explicitly retain the coefficient and transition-map hypotheses.

\begin{lemma}[AC-K29-irregular-separation: separation without ray idempotents]
\label{AC-K29-irregular-separation}
Let \(R\) be a commutative ring, let
\(G=H\rtimes\langle c\rangle\), with \(c^2=1\), and let
\(\theta:H\to R^\times\) be a character.  Put
\(\theta^c(h)=\theta(chc)\).  Suppose that for some \(h_0\in H\)
the element \(\theta(h_0)-\theta^c(h_0)\) is a unit.  Let \(J\) be
an \(R[G]\)-module on which, for every \(g\in G\), the characteristic
polynomial of \(\operatorname{Ind}_H^G\theta(g)\) annihilates the
operator \(g\).  There is an \(R\)-module \(M\) and an isomorphism
\[
 J\simeq \operatorname{Ind}_H^G\theta\otimes_R M.       \tag{CIS.1}
\]
In this isomorphism \(h\in H\) acts on the two coordinates through
\(\theta(h),\theta^c(h)\), and \(c\) exchanges the coordinates.
The construction uses the inverse of the single unit
\(\theta(h_0)-\theta^c(h_0)\), and no character-averaging idempotent.
\end{lemma}

\begin{proof}
Write \(U_h\) for the action of \(h\), and set
\(s_h=\theta(h)+\theta^c(h)\),
\(d_h=\theta(h)\theta^c(h)\).  The characteristic polynomials of
\(h\) and \(ch\) give
\[
 U_h^2-s_hU_h+d_h=0,\qquad (cU_h)^2=d_h.
\]
Consequently
\[
 cU_hc=d_hU_h^{-1}=s_h-U_h.                            \tag{CIS.2}
\]
For \(h,h'\in H\), apply this identity to \(hh'\) and also conjugate
the product \(U_hU_{h'}\) by \(c\).  The first computation gives
\(d_hd_{h'}U_{h'}^{-1}U_h^{-1}\), while the second gives
\(d_hd_{h'}U_h^{-1}U_{h'}^{-1}\).  All these operators are invertible,
so the operators \(U_h\) commute with one another.

Put \(a=\theta(h_0)\), \(b=\theta^c(h_0)\), and
\(P=(U_{h_0}-b)/(a-b)\).  Its defining polynomial gives
\(P^2=P\), and (CIS.2) gives \(cPc=1-P\).  Thus
\(J=M\oplus cM\), where \(M=PJ\).  Every \(U_h\) preserves these
two summands.  If \(B=U_h|_M\), its two characteristic-polynomial
identities, for \(h\) and \(h_0h\), are
\[
 B^2-(\theta(h)+\theta^c(h))B+\theta(h)\theta^c(h)=0,
\]
\[
 a^2B^2-a(a\theta(h)+b\theta^c(h))B
       +ab\theta(h)\theta^c(h)=0.
\]
Subtract \(a^2\) times the first identity from the second.  The result is
\(a\theta^c(h)(a-b)(B-\theta(h))=0\).  Its scalar coefficient is a
unit, so \(B=\theta(h)\).  Equation (CIS.2) gives the other character
on \(cM\).  This proves (CIS.1), including its stated normalization.
\end{proof}

\begin{lemma}[AC-K29-irregular-algebra: the canonical module on a CM branch]
\label{AC-K29-irregular-algebra}
Let \(T\) be a finite flat local \(A\)-algebra with residue field \(k\),
and write \(\omega_T=\operatorname{Hom}_A(T,A)\).  Consider a
surjection \(\phi:T\to A\) of one of the following kinds.

First, suppose that
\(T=A\times_k S\), where \(S\) is a finite flat local \(A\)-algebra
with residue field \(k\), that \(\phi\) is the first projection, and
that \(T\) has Cohen--Macaulay type three.  Second, for an integer
\(r\geq3\), suppose that
\[
 T=A\times_k\bigl(A[Z]/(Z^2-X^r)
                   \times_{A/(X^{r-1})}A\bigr),        \tag{CIS.3}
\]
where the map from the middle factor sends \(Z\) to zero, and that
\(\phi\) is either of the displayed \(A\)-component projections.
Then, as \(A\)-modules,
\[
 \omega_T\otimes_{T,\phi}A\simeq A\oplus k\oplus k.   \tag{CIS.4}
\]
\end{lemma}

\begin{proof}
For the first case choose an \(A\)-basis \(1,s_1,\ldots,s_d\) of
\(S\), with every \(s_i\) in its maximal ideal.  Such a basis exists
because \(1\notin XS\); subtracting constant residues from a basis
completing \(1\) gives the stated choice.  An \(A\)-basis of \(T\) is
\[
 1=(1,1),\quad q=(X,0),\quad t_i=(0,s_i).
\]
It satisfies \(q^2=Xq\), \(qt_i=0\).  If
\(s_is_j=c_{ij,0}+\sum_a c_{ij,a}s_a\), then
\(c_{ij,0}\in XA\) and
\[
 t_it_j=c_{ij,0}\,1-(c_{ij,0}/X)q+\sum_a c_{ij,a}t_a.
\]
Let \(1^*,q^*,t_i^*\) denote the dual basis of \(\omega_T\).
Tensoring along \(\phi\) means imposing \(q=X\) and \(t_i=0\).
The relations coming from \(q-X\) give
\(1^*=0\), \(Xt_i^*=0\).  The additional relations have no
\(q^*\)-component: for any functional \(v\),
\((t_i v)(q)=v(t_iq)=0\).  Therefore the quotient is
\(Aq^*\oplus k^a\) for an integer \(a\geq0\).  Reduction modulo
\(X\) identifies its fiber with
\(\omega_T/\mathfrak m_T\omega_T\).  The type-three hypothesis says
that the latter space has dimension three.  Hence \(a=2\).

For (CIS.3) use the following four-element \(A\)-basis:
\[
 1=(1,1,1),\quad q=(X,0,0),\quad a=(0,Z,0),
 \quad b=(0,0,X^{r-1}).
\]
Besides the identity, its multiplication table is
\[
 q^2=Xq,\quad qa=qb=ab=0,\quad
 a^2=X^r1-X^{r-1}q-Xb,\quad b^2=X^{r-1}b.             \tag{CIS.5}
\]
For the first component, impose \(q=X,a=b=0\) on the dual module.
The relations are exactly
\(1^*=0,\ Xa^*=0,\ Xb^*=0\); the class of \(q^*\) is free.
For the last component impose \(q=a=0,b=X^{r-1}\).
The relations are exactly
\(1^*=0,\ Xq^*=0,\ Xa^*=0\); the class of \(b^*\) is free.
These assertions follow by applying the transpose multiplication
matrices from (CIS.5).  For example, in the second case
\(a\cdot a^*=1^*\), \(a\cdot b^*=-Xa^*\), and
\(q\cdot q^*=1^*+Xq^*\).  All other relations follow from the three
displayed ones.  This proves (CIS.4) in both cases.
\end{proof}

\begin{proposition}[AC-K29-irregular-saturation: the actual new-level CM line]
\label{AC-K29-irregular-saturation}
Let \(g=\theta_\psi\) be a cuspidal weight-one theta series for an
imaginary quadratic field \(K\), where \(p\) splits, the conductor of
\(\psi\) is prime to \(p\), \(\psi\neq\psi^c\), and
\(\psi(\mathfrak p)=\psi(\overline{\mathfrak p})\).
Use the minimal tame level
\(N=|D_K|\operatorname N(\mathfrak f_\psi)\).
Let \(T\) be the completed characteristic-zero ordinary Hecke algebra
at its unique \(p\)-stabilization, after the strict coefficient
extension used in Betina--Dimitrov.  Thus the residue field is
\(\overline{\mathbb Q}_p\); equivalently one may work over a
coefficient field \(E\) to which the displayed ring and branch
descriptions descend.  We do not identify an arbitrary unsplit
\(E\)-local algebra with those fiber products.  Let \(\mathcal H\) be the completed
ordinary cohomology module, and choose either of the two CM maps
\(\phi:T\to A\).  Unramified rank-one coefficient twists are
understood in identifying the underlying modules.  Set
\(J=\mathcal H\otimes_{T,\phi}A\), and let
\(J^{\mathrm{tf}}\) be its quotient by \(A\)-torsion.

The image of Ohta's ordinary line gives an isomorphism onto the
corresponding rank-one character line in the free rank-two induced
module \(J^{\mathrm{tf}}\).  In particular that image remains nonzero
at \(X=0\); there is no positive power of \(X\) in this comparison.
More precisely, as \(A\)-modules with the global Galois action,
\[
 J\simeq\operatorname{Ind}_{G_K}^{G_{\mathbb Q}}\Theta
             \otimes_A(A\oplus k),\qquad
 J^{\mathrm{tf}}\simeq
             \operatorname{Ind}_{G_K}^{G_{\mathbb Q}}\Theta.       \tag{CIS.6}
\]
The ordinary line maps isomorphically onto its first character line
in the second isomorphism, after choosing that line's generator.
\end{proposition}

\begin{proof}
The primary inputs are Betina--Dimitrov, arXiv:1907.09422v3,
\S3.1 (the minimal level and (24)), Proposition 3.2, Theorem 4.8,
and Proposition 5.4; and Ohta's exact sequence and duality in the
precise form recalled in Betina--Maksoud--Pozzi, accepted author
version, \S7.1, equations (70)--(72), printed pp.~41--42.
The first source gives the two induced CM branches and exactly the
ring alternatives in the preceding lemma.  If both
anticyclotomic \(\mathscr L\)-invariants are nonzero, either CM
branch is an \(A\)-factor in a fiber product over \(k\), and
Proposition 5.4 computes the type to be three.  If one invariant
vanishes, Theorem 4.8(ii) gives (CIS.3), including its map
\(Z\mapsto0\).  No conjectural nonvanishing of both invariants is used.

Ohta's results give an exact sequence of underlying \(T\)-modules
\[
 0\longrightarrow T\longrightarrow\mathcal H
              \longrightarrow\omega_T\longrightarrow0.        \tag{CIS.7}
\]
After tensoring with \(A\), it is right exact.  The first map remains
injective: it is nonzero at the generic point of the CM branch,
where classical control identifies the ordinary line of the
rank-two representation, and its source \(A\) is torsionfree.
Thus (CIS.4) gives an exact sequence
\[
 0\longrightarrow A\xrightarrow{\iota}J
                 \longrightarrow A\oplus k^2\longrightarrow0.  \tag{CIS.8}
\]

The ordinary cohomology satisfies the characteristic-polynomial
relations of its Hecke pseudorepresentation.  They hold before
base change, by the Eichler--Shimura relations at unramified
Frobenius elements, their density, and continuity; equivalently
one may use the characteristic-polynomial identity on its
torsionfree, generically rank-two \(T\)-module.  They therefore
hold on \(J\), including its torsion.  On the CM branch they are
the characteristic polynomials of \(\operatorname{Ind}\Theta\).
Since \(\psi\neq\psi^c\), an element \(h_0\in G_K\) separates
the two characters by a unit of \(A\).  The separation lemma gives
\(J=M\oplus cM\), with \(M\) a finite \(A\)-module of rank one.

There is an inertia element at \(\mathfrak p\) for which
\(\Theta(s)-\Theta^c(s)\) is a unit times \(X\), and the ordinary
line has character \(\Theta\).  Its projection to \(cM\) is
therefore killed by \(X\); since it is the image of a map from
\(A\), this observation alone would not exclude a torsion
component.  To handle this point, use the global separation
projector \(P\) from (CIS.1).  Its restriction to the ordinary
line equals the identity at the generic point.  The endomorphism
\((1-P)\iota\) lands in \(J[X^\infty]\), and needs to be retained
until saturation is proved.

Write \(M=Ae\oplus C\), with \(C\) torsion, and write the image
of \(1\) under \(\iota\) as
\((X^a u e,0)+(t_1,t_2)\), where \(u\in A^\times\),
\(a\geq0\), and \(t_i\in C\).  The zero in the second free
coordinate follows from generic ordinarity.  Let \(b\) be the
length of \(C\).  The torsion length in the cokernel is
\(a+2b\), so (CIS.8) implies
\[
                         a+2b=2.                    \tag{CIS.9}
\]
If \(b=0\), then \(a=2\), but the cokernel is
\(A\oplus A/(X^2)\), contrary to (CIS.8).  Hence \(b=1\) and
\(a=0\).  In particular \(C\simeq k\), and the ordinary line
maps by the unit \(u\) onto the first free character line in
\(J^{\mathrm{tf}}\).  Its torsion coordinates disappear in that
quotient.  This proves (CIS.6) and the asserted saturation.
\end{proof}

\begin{remark}[Exact remaining ray-level interface]
The proposition is at the minimal tame level of each theta series.
It does not identify a theta specialization at an arbitrarily
raised tame level with that minimal-level module.  In a ray family
some characters have smaller conductor than the displayed modulus;
the resulting old vectors and their chosen tame \(U_q\)-eigenvalues
must be compared using the actual norm correspondences.
The printed integral patching theorem LLZ15, Proposition 5.2.5 and
Corollary 5.2.6, invokes the non-Eisenstein, \(p\)-distinguished
freeness theorem.  BL18, \S3.2.3, Propositions 3.7--3.10 retains
that hypothesis explicitly.  Neither theorem, as printed, supplies
this missing repeated-root ray-level comparison.

There is a useful exact descent principle once that comparison has
been proved.  If \(R=\mathcal O[\Delta][[X]]\) for a finite group
\(\Delta\), and a map \(F:J\to R^2\) is known to satisfy
\(F(J)[1/p]\subset XR[1/p]^2\), then
\(F(J)\subset XR^2\): reduce modulo \(X\) and use the
\(\mathcal O\)-torsionfreeness of the explicit group ring
\(R/XR=\mathcal O[\Delta]\).  This statement concerns that explicit
target ring, not an unspecified Hecke quotient.  It does not assert
that the hypothesis holds for the actual ray maps, nor that their
normalizations agree with the canonical ordinary differential.
\end{remark}
